\documentclass[10pt,a4paper]{amsart}
\usepackage[margin=19mm]{geometry}
\usepackage{amsmath,amssymb,mathtools}
\usepackage[scr=rsfs,cal=euler]{mathalfa}
\usepackage{xfrac}
\usepackage[unicode,hidelinks,bookmarks=false]{hyperref}
\newcommand{\ie}{i.e.\ }
\newcommand{\cf}{cf.\ }
\newcommand{\resp}{respectively\ }
\newcommand{\Subsection}{\subsection}
\providecommand{\nobreakdash}{\nobreak\hbox{-}}
\setlength{\emergencystretch}{2em}
\multlinegap0pt
\usepackage{enumerate}
\usepackage{amssymb}
\usepackage{mathtools}
\usepackage{tikz}
\newtheorem{theorem}{Theorem}
\newtheorem{lemma}{Lemma}
\def \Ad{\mathcal{A}_{k,l}}
\def\M{\mathcal{M}}
\title{Hartee-type heat equation associated to fractional anharmonic oscillator on weighted modulation spaces}
\author{Aparajita Dasgupta, Uttam Kumar Dolai}
\date{Source: arXiv:2606.11960v1 (CC BY 4.0)}
\begin{document}
\maketitle

\noindent\emph{Source and licence.} Hartee-type heat equation associated to fractional anharmonic oscillator on weighted modulation spaces. Version arXiv:2606.11960v1 (CC BY 4.0), distributed by arXiv under the Creative Commons Attribution 4.0 International licence, http://creativecommons.org/licenses/by/4.0/, as recorded in the arXiv licence metadata for this exact version and shown on the abstract page https://arxiv.org/abs/2606.11960v1. Full licence notice: \texttt{licenses/arxiv-2606.11960-CC-BY-4.0.txt}.\par\smallskip

\noindent\textit{Editorial scope.} Counted unit: the article's theorem welposedness2 (source0.tex lines 2821--2861). The article's complete original proof of it is delivered with it (source0.tex lines 2863--3208, one \texttt{proof} environment of 346 lines). No mathematics is written, repaired or re-derived: the statement and the proof are the article's own lines, in the article's own notation, and no retained line is substituted or re-flowed.  Retained uncounted as the source-local context the proof needs: lines 986--1031; lines 1552--1560; lines 1708--1733; lines 1736--1741; lines 2563--2617; lines 2633--2646.  Nothing was omitted from any retained span, so the package carries no omission record.  No retained line contains a citation, so the wrapper carries no bibliography and no bibliography entry.  No retained line contains a figure environment, an image-inclusion command or drawing code, so no image is delivered and the package declares no asset dependency.  Editorial formatting only, in addition: an AMS \texttt{amsart} preamble that restates the article's own declarations and macros used by the retained lines, namely the environments \texttt{theorem}, \texttt{lemma} on separate counters, together with the standard packages those lines need.  The retained environments print in this document as Theorem 1, Lemma 1 in order of appearance, whereas the article prints the same items with the numbering of its own full document.  The complete native source of the article as distributed in the arXiv e-print for this exact version is bundled unchanged as \texttt{sources/202589/source0.tex}.
\begin{theorem}
\label{BoundofH}
Let $\beta>0$, $0<p_{1},p_{2},q_{1},q_{2}\leq\infty$, $s_{1}>0$, and $s_{2}\in\mathbb{R}$. Define
\begin{align*}
\frac{1}{\widetilde{p}}
:=
\max\Bigl\{\frac{1}{p_{2}}-\frac{1}{p_{1}},\,0\Bigr\},
\qquad
\frac{1}{\widetilde{q}}
:=
\max\Bigl\{\frac{1}{q_{2}}-\frac{1}{q_{1}},\,0\Bigr\},
\end{align*}
and set
\begin{align}
\label{sigma}
\sigma
:=
\frac{d}{2\beta}
\left(
\frac{1}{k\,\widetilde{p}}
+
\frac{1}{l\,\widetilde{q}}
\right).
\end{align}
Then, for every $0<t\leq1$, the fractional anharmonic heat semigroup satisfies
\begin{equation}
\label{maininequality}
\bigl\|e^{-t\mathcal{A}_{k,l}^{\beta}}f\bigr\|_{\mathcal{M}^{p_{2},q_{2}}_{s_{2}}}
\le
C_{0}\,t^{-\sigma}\,
\|f\|_{\mathcal{M}^{p_{1},q_{1}}_{s_{1}}},
\end{equation}
where $C_{0}>0$ is a constant independent of both $t$ and $f$.

Moreover, in the case $s_{1}=s_{2}=s\ge0$, the estimate sharpens to
\begin{equation}
\label{secondconstant}
\bigl\|e^{-t\mathcal{A}_{k,l}^{\beta}}f\bigr\|_{\mathcal{M}^{p_{2},q_{2}}_{s}}
\le
\begin{cases}
C_{0}\,t^{-\sigma}\,\|f\|_{\mathcal{M}^{p_{1},q_{1}}_{s}}, & 0<t\le1,\\[0.3em]
C_{0}\,e^{-t\lambda_{0}^{\beta}}\,\|f\|_{\mathcal{M}^{p_{1},q_{1}}_{s}}, & t\ge1,
\end{cases}
\end{equation}
where $\lambda_{0}>0$ denotes the lowest eigenvalue of the anharmonic oscillator $\mathcal{A}_{k,l}$.
\end{theorem}

\begin{align}
\label{hartee1}
\begin{cases}
\partial_{t}u(t,x)+\Ad^{\beta}u(t,x)
=
\bigl(|x|^{-\gamma}*|u(t,x)|^{2}\bigr)u(t,x),\\
u(0,x)=u_{0}(x),
\end{cases}
\end{align}

\begin{theorem}
\label{welposedness1}
Let $0<\gamma<d$, $s>\frac{d}{q'}$, and
\[
1<p<\frac{d}{d-\gamma}.
\]
Assume that $r>0$ satisfies
\begin{align}
\label{r}
r>\frac{4\beta k}{2\beta k-d+\gamma}.
\end{align}
Then there exists $\varepsilon>0$ such that, whenever
\[
\|u_{0}\|_{\M^{p,q}_{s}}<\varepsilon,
\]
the nonlinear Hartree equation \eqref{hartee1} admits a unique global solution
\[
u\in L^{\infty}\bigl([0,\infty),\M^{p,q}_{s}\bigr)
   \cap L^{r}\bigl([0,\infty),\M^{p,q}_{s}\bigr).
\]
Moreover, the solution satisfies
\[
u\in C\bigl([0,\infty),\M^{p,q}_{s}\bigr)
   \cap L^{r}\bigl([0,\infty),\M^{p,q}_{s}\bigr).
\]
\end{theorem}

\begin{align}
\label{Duhamel}
\mathcal{J}(u)
:=e^{-t\Ad^{\beta}}u_{0}
+\int_{0}^{t}e^{-(t-\tau)\Ad^{\beta}}F(\tau,x)\,\mathrm{d}\tau,
\end{align}

\begin{lemma}
\label{Trlnerlma}
Let $0<\gamma<d$, $s\geq 0$, and let
\[
1\leq p,q,p_{0},q_{0}\leq \infty
\]
satisfy
\[
\frac{1}{p_{0}}
=
\frac{3}{p}
+\frac{\gamma}{d}
-1,
\qquad
\frac{1}{q_{0}}
=
3\Bigl(\frac{1}{q}-\frac{2}{3}\Bigr).
\]
Then, for all $f,g,h\in \M^{p,q}_{s}$, the Hartree-type trilinear operator obeys
\begin{align}
\label{Trlnr1}
\bigl\|
\bigl(|\cdot|^{-\gamma}*(f\overline{g})\bigr)h
\bigr\|_{\M^{p_{0},q_{0}}_{s}}
\lesssim
\|f\|_{\M^{p,q}_{s}}
\|g\|_{\M^{p,q}_{s}}
\|h\|_{\M^{p,q}_{s}}.
\end{align}
Moreover, the nonlinear mapping
\[
f \mapsto \bigl(|\cdot|^{-\gamma}*|f|^{2}\bigr)f
\]
is locally Lipschitz from $\M^{p,q}_{s}$ into $\M^{p_{0},q_{0}}_{s}$, in the sense that
\begin{align}
\label{Trlnr2}
&
\bigl\|
\bigl(|\cdot|^{-\gamma}*|f|^{2}\bigr)f
-
\bigl(|\cdot|^{-\gamma}*|g|^{2}\bigr)g
\bigr\|_{\M^{p_{0},q_{0}}_{s}}
\nonumber\\
&\lesssim
\Bigl(
\|f\|_{\M^{p,q}_{s}}^{2}
+
\|f\|_{\M^{p,q}_{s}}
\|g\|_{\M^{p,q}_{s}}
+
\|g\|_{\M^{p,q}_{s}}^{2}
\Bigr)
\|f-g\|_{\M^{p,q}_{s}}.
\end{align}
\end{lemma}

\begin{align}
\label{T101}
\frac{1}{p_{0}}
=
\frac{1}{p}
+
\frac{1}{p_{1}},
\qquad
1+\frac{1}{q_{0}}
=
\frac{1}{q}
+
\frac{1}{q_{1}}.
\end{align}

\begin{theorem}
\label{welposedness2}

Let $s\geq 0$, $0<\gamma<d$, and let
\[
1<p,q,r<\infty
\]
satisfy
\[
\frac{3d}{2d-\gamma}
<
p
<
\frac{3d}{d-\gamma},
\qquad
1<q<\frac{3}{2},
\qquad
r>
\frac{2lq'\beta}{lq'\beta-d}.
\]
Then there exists $\varepsilon>0$ such that, whenever
\[
\|u_{0}\|_{\M^{p,q}_{s}}<\varepsilon,
\]
the nonlinear Hartree equation \eqref{hartee1} admits a unique global solution
\[
u
\in
L^{\infty}\bigl([0,\infty),\M^{p,q}_{s}\bigr)
\cap
L^{r}\bigl([0,\infty),\M^{p,q}_{s}\bigr).
\]
Moreover, the solution satisfies
\[
u
\in
C\bigl([0,\infty),\M^{p,q}_{s}\bigr)
\cap
L^{r}\bigl([0,\infty),\M^{p,q}_{s}\bigr).
\]
\end{theorem}

\begin{proof}
As in the proof of Theorem \ref{welposedness1}, we denote by
\[
X
=
L^{\infty}\bigl([0,\infty),\M^{p,q}_{s}\bigr)
\cap
L^{r}\bigl([0,\infty),\M^{p,q}_{s}\bigr),
\]
equipped with the norm
\[
\|u\|_{X}
=
\|u\|_{L^{\infty}_{t}\M^{p,q}_{s}}
+
\|u\|_{L^{r}_{t}\M^{p,q}_{s}},
\]
and by $\mathcal{J}$ the integral operator defined in
\eqref{Duhamel}.

Proceeding exactly as in the proof of Theorem
\ref{welposedness1}, we obtain
\begin{align}
\label{awk}
\|e^{-t\Ad^{\beta}}u_{0}\|_{X}
\leq
C_{1}'\|u_{0}\|_{\M^{p,q}_{s}},
\end{align}
for some constant $C_{1}'>0$.

Next, using Theorem \ref{BoundofH} together with the choice of
exponents $p_{0}$ and $q_{0}$ from Lemma \ref{Trlnerlma}, we estimate
\begin{align}
\label{qew}
&
\Bigl\|
\int_{0}^{t}
e^{-(t-\tau)\Ad^{\beta}}
F(\tau,\cdot)\,\mathrm{d}\tau
\Bigr\|_{L^{r}_{t}\M^{p,q}_{s}}
\nonumber\\
&\lesssim
\Bigl\|
\int_{0}^{t}
C(t-\tau)
\|F(\tau,\cdot)\|_{\M^{p_{0},q_{0}}_{s}}
\,\mathrm{d}\tau
\Bigr\|_{L^{r}_{t}}
\nonumber\\
&\lesssim
\Bigl\|
C(\tau)
*
\|u(\tau,\cdot)\|_{\M^{p,q}_{s}}^{3}
\Bigr\|_{L^{r}_{t}},
\end{align}
where
\[
C(\tau)
=
\begin{cases}
\tau^{-\sigma},
& 0<\tau\leq 1,
\\[0.3em]
e^{-\tau\lambda_{0}^{\beta}},
& \tau\geq 1.
\end{cases}
\]

The exponent $\sigma$ is given by
\begin{align}
\label{Sigma1}
\sigma
=
\frac{d}{2\beta}
\left(
\frac{1}{k\widetilde{p}}
+
\frac{1}{l\widetilde{q}}
\right).
\end{align}

Using \eqref{T101}, we compute
\[
\frac{1}{\widetilde{p}}
=
\max\Bigl\{
\frac{1}{p}-\frac{1}{p_{0}},
0
\Bigr\}
=
0,
\]
and
\[
\frac{1}{\widetilde{q}}
=
\max\Bigl\{
\frac{1}{q}-\frac{1}{q_{0}},
0
\Bigr\}
=
1-\frac{1}{q_{1}}
=
2-\frac{2}{q}
=
\frac{2}{q'}.
\]
Consequently, \eqref{Sigma1} reduces to
\begin{align}
\label{sigma2}
\sigma
=
\frac{d}{lq'\beta}.
\end{align}

Applying Young's inequality to \eqref{qew}, we obtain
\begin{align}
\label{qwee}
\Bigl\|
\int_{0}^{t}
e^{-(t-\tau)\Ad^{\beta}}
F(\tau,\cdot)\,\mathrm{d}\tau
\Bigr\|_{L^{r}_{t}\M^{p,q}_{s}}
\lesssim
\|C(t)\|_{L^{r_{2}}_{t}}
\Bigl\|
\|u\|_{\M^{p,q}_{s}}^{3}
\Bigr\|_{L^{r/3}_{t}},
\end{align}
where $r_{2}>1$ is chosen so that
\[
1+\frac{1}{r}
=
\frac{1}{r_{2}}
+
\frac{3}{r}.
\]

As in the proof of Theorem~\ref{welposedness1}, the quantity
\[
\|C(t)\|_{L^{r_{2}}_{t}}
\]
is finite whenever
\[
\sigma r_{2}<1.
\]
Since \(\frac{1}{r_{2}}=1-\frac{2}{r}\), this condition is equivalent to
\begin{equation}
\label{qwe1}
\sigma
<
1-\frac{2}{r}.
\end{equation}
It remains to verify that \eqref{qwe1} follows from the assumption on \(r\). Indeed, the condition
\[
r>
\frac{2lq'\beta}{lq'\beta-d}
\]
implies that
\[
\frac{2}{r}
<
1-\frac{d}{lq'\beta},
\]
and hence,
\[
\frac{d}{lq'\beta}
<
1-\frac{2}{r}.
\]
Invoking \eqref{sigma2}, we immediately obtain
\[
\sigma
<
1-\frac{2}{r},
\]
which is precisely condition \eqref{qwe1}.

Therefore, from \eqref{qwee}, we conclude that
\begin{align}
\label{asd}
\Bigl\|
\int_{0}^{t}
e^{-(t-\tau)\Ad^{\beta}}
F(\tau,\cdot)\,\mathrm{d}\tau
\Bigr\|_{L^{r}_{t}\M^{p,q}_{s}}
\lesssim
\|u\|_{L^{r}_{t}\M^{p,q}_{s}}^{3}.
\end{align}


Using the same argument as in the proof of Theorem
\ref{welposedness1}, together with the trilinear estimate
\eqref{Trlnr1}, we obtain
\begin{align}
\label{asd1}
\Bigl\|
\int_{0}^{t}
e^{-(t-\tau)\Ad^{\beta}}
F(\tau,\cdot)\,\mathrm{d}\tau
\Bigr\|_{L^{\infty}_{t}([0,\infty),\M^{p,q}_{s}}
\lesssim
\|u\|^{3}_{L^{\infty}_{t}([0,\infty),\M^{p,q}_{s}}.
\end{align}

Combining \eqref{asd} and \eqref{asd1}, we deduce that
\begin{align}
\label{awk1}
\Bigl\|
\int_{0}^{t}
e^{-(t-\tau)\Ad^{\beta}}
F(\tau,\cdot)\,\mathrm{d}\tau
\Bigr\|_{X}
\leq
C_{2}'\|u\|_{X}^{3},
\end{align}
for some constant $C_{2}'>0$.

Consequently, from \eqref{awk} and \eqref{awk1}, we arrive at
\begin{align}
\label{awk2}
\|\mathcal{J}(u)\|_{X}
\leq
C_{3}'
\Bigl(
\|u_{0}\|_{\M^{p,q}_{s}}
+
\|u\|_{X}^{3}
\Bigr),
\end{align}
for some constant $C_{3}'>0$.

Choose $\varepsilon>0$ such that
\[
\varepsilon^{2}
<
\frac{1}{2C_{3}'},
\]
and assume that the initial data satisfies
\[
\|u_{0}\|_{\M^{p,q}_{s}}
<
\frac{\varepsilon}{2C_{3}'}.
\]
Define
\[
B_{\varepsilon}
=
\Bigl\{
u\in X
:
\|u\|_{X}\leq \varepsilon
\Bigr\},
\]
which is a closed ball centered at the origin in $X$.
Then, by \eqref{awk2}, the operator $\mathcal{J}$ maps
$B_{\varepsilon}$ into itself.

{We next show that, for a sufficiently small choice of $\varepsilon>0$, the operator $\mathcal{J}$ is a contraction on $B_{\varepsilon}$.} Let
\[
F(t,x)
=
\bigl(|x|^{-\gamma}*|u(t,x)|^{2}\bigr)u(t,x),
\]
and
\[
G(t,x)
=
\bigl(|x|^{-\gamma}*|v(t,x)|^{2}\bigr)v(t,x).
\]
Using the estimate \eqref{Trlnr2}, together with Young's and
H\"older's inequalities, and arguing exactly as before, we obtain
\begin{align*}
&
\Bigl\|
\int_{0}^{t}
e^{-(t-\tau)\Ad^{\beta}}
\bigl(F(\tau,\cdot)-G(\tau,\cdot)\bigr)
\,\mathrm{d}\tau
\Bigr\|_{L^{r}_{t}\M^{p,q}_{s}}
\\
&\lesssim
\Bigl(
\|u\|_{L^{r}_{t}\M^{p,q}_{s}}^{2}
+
\|u\|_{L^{r}_{t}\M^{p,q}_{s}}
\|v\|_{L^{r}_{t}\M^{p,q}_{s}}
+
\|v\|_{L^{r}_{t}\M^{p,q}_{s}}^{2}
\Bigr)
\|u-v\|_{L^{r}_{t}\M^{p,q}_{s}},
\end{align*}
and similarly,
\begin{align*}
&
\Bigl\|
\int_{0}^{t}
e^{-(t-\tau)\Ad^{\beta}}
\bigl(F(\tau,\cdot)-G(\tau,\cdot)\bigr)
\,\mathrm{d}\tau
\Bigr\|_{L^{\infty}_{t}\M^{p,q}_{s}}
\\
&\lesssim
\Bigl(
\|u\|_{L^{\infty}_{t}\M^{p,q}_{s}}^{2}
+
\|u\|_{L^{\infty}_{t}\M^{p,q}_{s}}
\|v\|_{L^{\infty}_{t}\M^{p,q}_{s}}
+
\|v\|_{L^{\infty}_{t}\M^{p,q}_{s}}^{2}
\Bigr)
\|u-v\|_{L^{\infty}_{t}\M^{p,q}_{s}}.
\end{align*}

Since $u,v\in B_{\varepsilon}$, choosing $\varepsilon>0$
sufficiently small and combining the above two estimates, we obtain
\[
\|\mathcal{J}(u)-\mathcal{J}(v)\|_{X}
\leq
\frac{1}{2}
\|u-v\|_{X}.
\]
Hence $\mathcal{J}$ is a contraction on $B_{\varepsilon}$.

Therefore, by Banach's contraction mapping principle,
$\mathcal{J}$ admits a unique fixed point in $B_{\varepsilon}$,
which yields a unique global solution to \eqref{hartee1} in
\[
L^{\infty}([0,\infty),\M^{p,q}_{s})
\cap
L^{r}([0,\infty),\M^{p,q}_{s}),
\]
for all exponents $p,q,r,s$ satisfying the assumptions of the theorem.

Finally, arguing exactly as in the proof of Theorem
\ref{welposedness1}, we conclude that the solution moreover satisfies
\[
u
\in
C\bigl([0,\infty),\M^{p,q}_{s}\bigr)
\cap
L^{r}([0,\infty),\M^{p,q}_{s}).
\]
This completes the proof.
 \end{proof}

\end{document}
